% !TeX root = ../../Book1.tex
% 实读来源：Rudin1991Functional，13.34--13.35，印刷 pp.375--379；
% 13.38，印刷 pp.382--385。大学公开扫描的零基 PDF 页序为印刷页码加 16。
% 目视核对 pp.379、382；正文自证酉群情形的平滑平均、两侧 Laplace 公式与唯一性，
% 不调用该书 13.38 所采用的 Gelfand 正规半群构造。
\section{强连续酉群与 Stone 定理}
\label{b1:sec:2908}

谱积分可以把实变量函数 \(e^{\mathrm it\lambda}\) 变成一个酉算子。
当时间 \(t\) 改变时，这些算子满足群律；反过来，
一个连续而且保持内积的演化，是否一定来自某个自伴算子？
Stone 定理给出肯定回答。这里的连续是逐个初值的强连续，
零时导数则只对一个稠密子空间中的初值存在。
因此，对应关系必须同时确定算子的作用和定义域。

\ProseParagraphBreak

本节固定
\[
 U(t)=e^{\mathrm itA},\qquad
 Gx=\lim_{h\to0}\frac{U(h)x-x}{h}=\mathrm iAx.
\]
其中 \(G\) 是由差商定义的生成元，\(A=-\mathrm iG\) 是要证明自伴的算子。
关于生成元的基本性质与 Laplace 公式，可对照 Rudin 的
《Functional Analysis》\cite[Theorem~13.35, pp.~376--379; Theorem~13.38, pp.~382--385]{Rudin1991Functional}。
下面在酉群的情形逐步完成证明。

\subsection{强连续酉群}
\label{b1:subsec:290801}

按\cref{b1:def:29-unitary-equivalence}，酉算子是满射线性等距映射。
在同一个 Hilbert 空间上，这等价于 \(U^*U=UU^*=I\)。
现在用时间参数把这些保持内积的算子组织起来。

\begin{definition}[强连续酉群]
\label{b1:def:29-unitary-group}
设 \(H\) 是复 Hilbert 空间。若一族有界线性算子
\((U(t))_{t\in\mathbb R}\subset\mathcal B(H)\) 满足
\[
 U(0)=I,\qquad U(t+s)=U(t)U(s)\quad(s,t\in\mathbb R),
\]
每个 \(U(t)\) 都是酉算子，并且对每个 \(x\in H\)，
\[
 \lim_{t\to0}\|U(t)x-x\|=0,
\]
则称 \((U(t))_{t\in\mathbb R}\) 为 \(H\) 上的
\term[qiang-lianxu-you-qun]{强连续酉群}[strongly continuous unitary group]。
\end{definition}

群律和酉性给出
\[
 U(t)^*=U(-t),\qquad \|U(t)x\|=\|x\|.
\]
而且
\[
 \|U(t+h)x-U(t)x\|=\|U(h)x-x\|,
\]
所以零时的强连续自动给出每条轨道 \(t\mapsto U(t)x\) 在全部实轴上的连续性。
这没有要求算子范数 \(\|U(h)-I\|\) 趋于零。

\ProseParagraphBreak

例如，在复空间 \(\ell^2\) 上令
\[
 U(t)(x_n)_{n\geq1}=(e^{\mathrm i n t}x_n)_{n\geq1}.
\]
群律和酉性可以逐坐标核对。对固定 \(x=(x_n)\in\ell^2\)，有
\[
 \|(U(h)-I)x\|_2^2
 \leq\sum_{n=1}^N|e^{\mathrm i n h}-1|^2\cdot|x_n|^2
       +4\sum_{n>N}|x_n|^2.
\]
先选 \(N\) 控制尾和，再令 \(h\to0\) 控制有限和，便得强连续。
但令 \(h_n=\pi/n\)，并以第 \(n\) 个单位向量测试，便有
\[
 \|U(h_n)-I\|\geq\|(U(h_n)-I)e_n\|_2=2.
\]
另一方面，酉性给出上界 \(2\)，故这些范数恰等于 \(2\)。
强连续允许高频坐标随时间尺度改变，不能把固定初值的极限改成单位球上的一致极限。

\ProseParagraphBreak

证明生成元的性质需要积分连续的 \(H\) 值轨道。
这里所需积分可以直接用 Riemann 和建立。
若 \(v:[a,b]\to H\) 连续，则它一致连续；
把两个足够细的带标记分割细化到共同分割后，
相应 Riemann 和之差的范数由各小区间长度乘以 \(v\) 的连续模控制。
这些和因而构成 Cauchy 族，\(H\) 的完备性给出唯一极限，记为
\(\int_a^b v(s)\,\mathrm ds\)。从有限和取极限可得
\begin{equation}\label{b1:eq:29-vector-integral-bound}
 \left\|\int_a^b v(s)\,\mathrm ds\right\|
 \leq\int_a^b\|v(s)\|\,\mathrm ds.
\end{equation}
有界线性算子可与这个积分交换，内积对任一固定向量的测试也可与积分交换。
约定 \(\int_b^a=-\int_a^b\)，则对连续 \(v\)，
\[
 \left\|\frac1h\int_t^{t+h}v(s)\,\mathrm ds-v(t)\right\|
 \leq\sup_{|s-t|\leq|h|}\|v(s)-v(t)\|\longrightarrow0.
\]
这就是所需的强微积分基本公式。
若 \(F:[a,b]\to H\) 连续可微，对任一 \(y\in H\) 将标量微积分基本定理
用于 \(\langle F(s),y\rangle\)，便得
\(F(b)-F(a)=\int_a^bF'(s)\,\mathrm ds\)。
以下有限区间积分均取这个意义，无穷区间积分则取有限区间积分的范数极限。

\subsection{生成元、符号与定义域}
\label{b1:subsec:290802}

\begin{definition}[酉群的生成元]
\label{b1:def:29-group-generator}
设 \(H\) 是复 Hilbert 空间，\((U(t))_{t\in\mathbb R}\) 是其上的强连续酉群。
定义
\begin{equation}\label{b1:eq:29-generator-domain}
 \begin{split}
 D(G)&=\left\{x\in H:
       \lim_{h\to0}\frac{U(h)x-x}{h}\text{ 在 }H\text{ 中存在}\right\},\\
 Gx&=\lim_{h\to0}\frac{U(h)x-x}{h}\qquad(x\in D(G)).
 \end{split}
\end{equation}
线性算子 \(G:D(G)\subset H\to H\) 称为该群的
\term[shengchengyuan]{生成元}[generator]。
极限取实参数 \(h\) 的双侧极限，收敛均指 \(H\) 的范数收敛。
\end{definition}

差商的线性表明 \(D(G)\) 是线性子空间，\(G\) 在其中线性。
对群而言，只要求右侧差商收敛也会得到同一个定义域：若右极限为 \(y\)，
则对 \(h>0\)，
\[
 \frac{U(-h)x-x}{-h}
 =U(-h)\frac{U(h)x-x}{h}\longrightarrow y.
\]
最后一步可把右端与 \(y\) 的差拆成
\(U(-h)((U(h)x-x)/h-y)+(U(-h)y-y)\)，
再用酉性和强连续性分别控制两项。

\begin{proposition}[定义域不变与轨道积分式]
\label{b1:prop:29-group-orbits}
设 \(H\) 是复 Hilbert 空间，\((U(t))_{t\in\mathbb R}\) 是其上的强连续酉群，
\(G\) 为其生成元。对每个 \(t\in\mathbb R\)，有
\[
 U(t)D(G)=D(G),\qquad GU(t)x=U(t)Gx\quad(x\in D(G)).
\]
对每个 \(x\in D(G)\)，轨道是 \(\mathbb R\) 上连续可微的 \(H\) 值函数，且
\begin{equation}\label{b1:eq:29-group-orbit-derivative}
 \frac{\mathrm d}{\mathrm dt}U(t)x=U(t)Gx=GU(t)x,
 \qquad
 U(t)x-x=\int_0^tU(s)Gx\,\mathrm ds.
\end{equation}
\end{proposition}

\begin{proof}
若 \(x\in D(G)\)，群律给出
\[
 \frac{U(h)U(t)x-U(t)x}{h}
 =U(t)\frac{U(h)x-x}{h}\longrightarrow U(t)Gx.
\]
故 \(U(t)x\in D(G)\)，并得到交换公式。
再把时间换成 \(-t\)，可得定义域之间反向的包含，从而是等式。
同一差商也给出 \(t\) 时刻的轨道导数。
由于 \(s\mapsto U(s)Gx\) 强连续，导数连续，
上一小节的微积分基本公式便给出积分恒等式，包括负时间的情形。
\end{proof}

\begin{proposition}[稠密性、闭性与内积关系]
\label{b1:prop:29-generator-closed}
设 \(H\) 是复 Hilbert 空间，\((U(t))_{t\in\mathbb R}\) 是其上的强连续酉群，
\(G:D(G)\subset H\to H\) 为其生成元。
则 \(D(G)\) 在 \(H\) 中稠密，\(G\) 是闭算子，并且
\begin{equation}\label{b1:eq:29-generator-skew-symmetry}
 \langle Gx,y\rangle+\langle x,Gy\rangle=0
 \qquad(x,y\in D(G)).
\end{equation}
因此 \(A=-\mathrm iG\)、\(D(A)=D(G)\) 是稠密定义的闭对称算子。
\end{proposition}

\begin{proof}
先以平滑的时间平均逼近任意 \(x\in H\)。
取非负 \(\varphi\in C_c^1(\mathbb R)\)，支集包含于 \([-1,1]\)，且
\(\int_{\mathbb R}\varphi(s)\,\mathrm ds=1\)；
例如把 \((1-s^2)^2\) 在 \([-1,1]\) 外延为零，再作归一化。
对 \(\varepsilon>0\)，令
\[
 \varphi_\varepsilon(s)=\varepsilon^{-1}\varphi(s/\varepsilon),\qquad
 x_\varepsilon=\int_{\mathbb R}\varphi_\varepsilon(s)U(s)x\,\mathrm ds.
\]
这些积分的被积函数连续且具有紧支集。由积分估计，
\[
 \|x_\varepsilon-x\|
 \leq\int\varphi_\varepsilon(s)\|U(s)x-x\|\,\mathrm ds
 \leq\sup_{|s|\leq\varepsilon}\|U(s)x-x\|\longrightarrow0.
\]
还需证明每个 \(x_\varepsilon\) 确实属于生成元域。
将有界算子 \(U(h)\) 移入积分，再以 \(s\) 代替原变量加 \(h\)，得到
\[
 \frac{U(h)x_\varepsilon-x_\varepsilon}{h}
 =\int_{\mathbb R}
   \frac{\varphi_\varepsilon(s-h)-\varphi_\varepsilon(s)}{h}
   U(s)x\,\mathrm ds.
\]
标量差商等于
\(-\int_0^1\varphi_\varepsilon'(s-\theta h)\,\mathrm d\theta\)。
由于 \(\varphi_\varepsilon'\) 一致连续，它在 \(h\to0\) 时一致趋于
\(-\varphi_\varepsilon'(s)\)；当 \(|h|<1\) 时，所有被积函数又有共同的紧支集。
用 \(\|U(s)x\|=\|x\|\) 及积分估计，可将极限移入积分，得到
\[
 x_\varepsilon\in D(G),\qquad
 Gx_\varepsilon=-\int_{\mathbb R}\varphi_\varepsilon'(s)U(s)x\,\mathrm ds.
\]
故 \(D(G)\) 稠密。

为证闭性，设 \(x_n\in D(G)\)、\(x_n\to x\) 且 \(Gx_n\to y\)。
由\eqref{b1:eq:29-group-orbit-derivative}，
\[
 U(t)x_n-x_n=\int_0^tU(s)Gx_n\,\mathrm ds.
\]
左端趋于 \(U(t)x-x\)，右端与 \(\int_0^tU(s)y\,\mathrm ds\) 的差的范数
不超过 \(|t|\cdot\|Gx_n-y\|\)。因此对每个 \(t\in\mathbb R\)，
\[
 U(t)x-x=\int_0^tU(s)y\,\mathrm ds.
\]
除以非零 \(t\) 后令 \(t\to0\)，积分平均由强连续性趋于 \(y\)。
于是 \(x\in D(G)\)、\(Gx=y\)，这就证明了闭性。

最后，对 \(x,y\in D(G)\)，酉性给出
\(\langle U(t)x,U(t)y\rangle=\langle x,y\rangle\)。
在零时求导，便得\eqref{b1:eq:29-generator-skew-symmetry}。
本书内积对第一变量线性，所以
\[
 \langle-\mathrm iGx,y\rangle
 =-\mathrm i\langle Gx,y\rangle
 =\mathrm i\langle x,Gy\rangle
 =\langle x,-\mathrm iGy\rangle.
\]
这证明 \(A=-\mathrm iG\) 对称；其稠密性和闭性由 \(G\) 的相应性质给出。
\end{proof}

对称性尚未说明 \(A\) 自伴。
这里仍缺少伴随定义域的辨认，最后一小节将借助两个非实参数处的满射性补齐它。

\subsection{从自伴算子到酉群}
\label{b1:subsec:290803}

先从自伴算子构造演化。在这个方向，前面建立的谱演算同时给出
算子、连续性以及准确的可微初值范围。

\begin{proposition}[自伴算子的指数群]
\label{b1:prop:29-spectral-unitary-group}
设 \(H\) 是复 Hilbert 空间，\(A:D(A)\subset H\to H\) 是自伴算子，
\(E\) 是它在 \(\mathbb R\) 上的投影值谱测度。对 \(t\in\mathbb R\)，定义
\begin{equation}\label{b1:eq:29-unitary-spectral-integral}
 U_A(t)=e^{\mathrm itA}
       =\int_{\mathbb R}e^{\mathrm it\lambda}\,\mathrm dE(\lambda).
\end{equation}
则 \((U_A(t))_{t\in\mathbb R}\) 是 \(H\) 上的强连续酉群，
它的生成元恰为
\[
 D(G_A)=D(A),\qquad G_Ax=\mathrm iAx\quad(x\in D(A)).
\]
对每个 \(t\in\mathbb R\)，有
\[
 U_A(t)D(A)=D(A),\qquad AU_A(t)x=U_A(t)Ax\quad(x\in D(A)),
\]
且对每个 \(x\in D(A)\)，
\[
 \frac{\mathrm d}{\mathrm dt}U_A(t)x
 =\mathrm iAU_A(t)x=\mathrm iU_A(t)Ax
\]
是连续的强导数。
\end{proposition}

\begin{proof}
由无界自伴谱定理\cref{b1:thm:29-unbounded-spectral}及相应函数演算，
\(|e^{\mathrm it\lambda}|=1\) 保证 \(U_A(t)\) 是定义在全部 \(H\) 上的有界算子。
有界 Borel 演算的乘法与伴随规则给出
\[
 U_A(0)=I,\quad U_A(t+s)=U_A(t)U_A(s),\quad
 U_A(t)^*=U_A(-t)=U_A(t)^{-1}.
\]
对 \(x\in H\)，记 \(\mu_x(B)=\langle E(B)x,x\rangle\)。
谱范数公式给出
\[
 \|U_A(h)x-x\|^2
 =\int_{\mathbb R}|e^{\mathrm ih\lambda}-1|^2\,\mathrm d\mu_x(\lambda).
\]
被积函数逐点趋于零，且不超过 \(4\)，而 \(\mu_x(\mathbb R)=\|x\|^2\)。
支配收敛定理于是给出零时强连续，群律再给出全部时间的强连续。

若 \(x\in D(A)\)，则 \(\int\lambda^2\,\mathrm d\mu_x<\infty\)。
对实数 \(\lambda\) 及非零实数 \(h\)，标量恒等式
\[
 \frac{e^{\mathrm ih\lambda}-1}{h}
 =\mathrm i\lambda\int_0^1e^{\mathrm i\theta h\lambda}\,\mathrm d\theta
\]
给出其模不超过 \(|\lambda|\)，而点态极限为 \(\mathrm i\lambda\)。
因而
\[
 \begin{split}
 &\left\|\frac{U_A(h)x-x}{h}-\mathrm iAx\right\|^2\\
 &\qquad=\int_{\mathbb R}
  \left|\frac{e^{\mathrm ih\lambda}-1}{h}-\mathrm i\lambda\right|^2
  \,\mathrm d\mu_x(\lambda)\longrightarrow0,
 \end{split}
\]
这里误差平方由可积函数 \(4\lambda^2\) 控制。
这证明 \(D(A)\subseteq D(G_A)\)，以及在该域上的作用公式。

反过来，设 \(x\in D(G_A)\)。差商收敛使其范数在充分小的时间内有界。
取 \(h_n=1/n\)，由 Fatou 引理及谱范数公式，
\[
 \begin{split}
 \int_{\mathbb R}\lambda^2\,\mathrm d\mu_x(\lambda)
 &\leq\liminf_{n\to\infty}
  \int_{\mathbb R}
   \left|\frac{e^{\mathrm ih_n\lambda}-1}{h_n}\right|^2
   \,\mathrm d\mu_x(\lambda)\\
 &=\liminf_{n\to\infty}
  \left\|\frac{U_A(h_n)x-x}{h_n}\right\|^2<\infty.
 \end{split}
\]
故 \(x\in D(A)\)。于是定义域相等，而不只是在预先给定的域上导数相同。
上述必要性论证甚至只需要右侧差商存在。

还可直接从谱测度看出定义域为什么不变。
投影 \(E(B)\) 与 \(U_A(t)\) 交换，酉性给出
\[
 \mu_{U_A(t)x}(B)
 =\|E(B)U_A(t)x\|^2
 =\|U_A(t)E(B)x\|^2=\mu_x(B).
\]
因此 \(\int\lambda^2\,\mathrm d\mu_{U_A(t)x}
=\int\lambda^2\,\mathrm d\mu_x\)，这给出双向的域不变性。
若 \(x\in D(A)\)，令
\(A_N=\int\lambda\cdot\mathbf1_{[-N,N]}(\lambda)\,\mathrm dE(\lambda)\)。
有界乘法规则给出 \(A_NU_A(t)x=U_A(t)A_Nx\)；
两端分别趋于 \(AU_A(t)x\) 和 \(U_A(t)Ax\)，故无界算子的交换公式成立。
最后应用\cref{b1:prop:29-group-orbits}，便得到全部时间的强导数。
\end{proof}

把实乘法算子定理\cref{b1:thm:29-real-multiplication}用于 \(\mathbb N\) 上的计数测度，
开头的对角群便对应于自伴算子 \(Ax=(nx_n)\)，其最大定义域为
\[
 D(A)=\left\{x\in\ell^2:\sum_{n=1}^{\infty}n^2|x_n|^2<\infty\right\},
 \qquad Gx=(\mathrm i n x_n).
\]
这也可以由逐坐标差商与平方和的 Fatou 引理直接核对。
例如 \(x=(1/n)_{n\geq1}\) 属于 \(\ell^2\)，却不属于 \(D(A)\)。
其轨道处处连续，但在任何时间都没有强导数：
若某时刻差商收敛，左乘该时刻演化的逆，就会得到零时差商收敛，
与定义域的判别矛盾。

\subsection{Stone 对应与唯一性}
\label{b1:subsec:290804}

给定一个强连续酉群，\cref{b1:prop:29-generator-closed}已证明 \(-\mathrm iG\) 闭且对称。
要应用自伴满射判据，需要证明两个相反方向的非实参数都没有可解性缺口。
正、负时间轨道的 Laplace 积分恰好提供这两个逆。

\begin{lemma}[两侧 Laplace 公式]
\label{b1:lem:29-group-laplace}
设 \(H\) 是复 Hilbert 空间，\((U(t))_{t\in\mathbb R}\) 是其上的强连续酉群，
\(G\) 是其生成元。对 \(\lambda>0\)，定义
\[
 \begin{split}
 B_\lambda^+x&=\int_0^\infty e^{-\lambda s}U(s)x\,\mathrm ds,\\
 B_\lambda^-x&=\int_{-\infty}^0e^{\lambda s}U(s)x\,\mathrm ds
              =\int_0^\infty e^{-\lambda s}U(-s)x\,\mathrm ds
              \qquad(x\in H).
 \end{split}
\]
则 \(B_\lambda^\pm\in\mathcal B(H)\)、\(\|B_\lambda^\pm\|\leq\lambda^{-1}\)，且
\[
 (\lambda I-G)^{-1}=B_\lambda^+,\qquad
 (\lambda I+G)^{-1}=B_\lambda^-.
\]
这里 \(\lambda I\mp G:D(G)\to H\) 均为双射，
各逆的值域恰为 \(D(G)\)。
\end{lemma}

\begin{proof}
先逐个固定 \(x\in H\) 作积分。
酉性和\eqref{b1:eq:29-vector-integral-bound}给出
\[
 \left\|\int_T^R e^{-\lambda s}U(\pm s)x\,\mathrm ds\right\|
 \leq\|x\|\int_T^R e^{-\lambda s}\,\mathrm ds
 \leq\frac{e^{-\lambda T}}{\lambda}\|x\|
 \qquad(0\leq T<R).
\]
因此无穷积分按范数收敛；其线性及范数估计说明它们定义了有界算子。
这个构造没有把 \(s\mapsto U(s)\) 当作按算子范数连续的函数来积分。

令 \(b=B_\lambda^+x\)。将 \(U(h)\) 移入有限区间积分、换元，再取无穷积分极限，得到
\[
 U(h)b=e^{\lambda h}
       \left(b-\int_0^h e^{-\lambda r}U(r)x\,\mathrm dr\right).
\]
此式对正、负 \(h\) 均成立，负上限使用有向积分的约定。
于是
\[
 \frac{U(h)b-b}{h}
 =\frac{e^{\lambda h}-1}{h}b
  -\frac{e^{\lambda h}}{h}
   \int_0^h e^{-\lambda r}U(r)x\,\mathrm dr
 \longrightarrow\lambda b-x.
\]
故 \(b\in D(G)\)，且 \((\lambda I-G)B_\lambda^+x=x\)。
反向时间族 \(V(t)=U(-t)\) 仍为强连续酉群，
由双侧差商定义，它的生成元是 \(-G\)，定义域仍为 \(D(G)\)。
把刚才的计算用于 \(V\)，便得
\((\lambda I+G)B_\lambda^-x=x\)。这已经证明两个算子满射。

为证单射，对 \(y\in D(G)\)，由\eqref{b1:eq:29-generator-skew-symmetry}可知
\(\operatorname{Re}\langle Gy,y\rangle=0\)，从而
\[
 \|(\lambda I\mp G)y\|^2
 =\lambda^2\|y\|^2+\|Gy\|^2
 \geq\lambda^2\|y\|^2.
\]
所以两者均为单射，已经构造的右逆就是双侧逆。
具体地，对 \(y\in D(G)\)，将
\(B_\lambda^\pm(\lambda I\mp G)y-y\) 代入 \(\lambda I\mp G\)，结果为零，
单射性便迫使原向量为零。这也证明逆的值域恰为整个 \(D(G)\)。
\end{proof}

下面的\term[stone-dingli]{Stone 定理}[Stone's theorem]把自伴算子与强连续酉群一一对应起来。

\begin{theorem}[Stone]
\label{b1:thm:29-stone}
设 \(H\) 是复 Hilbert 空间。对 \(H\) 上的每个强连续酉群
\((U(t))_{t\in\mathbb R}\)，存在唯一自伴算子
\(A:D(A)\subset H\to H\)，使
\[
 U(t)=e^{\mathrm itA}\qquad(t\in\mathbb R),
\]
其中右端按谱演算定义。若 \(G\) 是该群的生成元，则准确地有
\[
 D(A)=D(G),\qquad A=-\mathrm iG,\qquad G^*=-G.
\]
反过来，每个自伴算子通过 \(e^{\mathrm itA}\) 给出一个强连续酉群，
其生成元为 \(\mathrm iA\)，定义域为 \(D(A)\)。
具有相同生成元的两个强连续酉群必然相同；这里相同包括定义域相同。
\end{theorem}

\begin{proof}
给定强连续酉群，令 \(A=-\mathrm iG\)、\(D(A)=D(G)\)。
\cref{b1:prop:29-generator-closed}表明它稠密定义、闭且对称。
由\cref{b1:lem:29-group-laplace}在 \(\lambda=1\) 时的满射性，
\[
 A+\mathrm iI=\mathrm i(I-G),\qquad
 A-\mathrm iI=-\mathrm i(I+G)
\]
都有值域 \(H\)。自伴满射判据\cref{b1:thm:29-self-adjoint-criterion}
于是给出 \(A=A^*\)。由于 \(G=\mathrm iA\)，伴随定义也给出
\(D(G^*)=D(A)\)、\(G^*=-\mathrm iA=-G\)。

现在由\cref{b1:prop:29-spectral-unitary-group}构造
\(V(t)=e^{\mathrm itA}\)。这个群的生成元是 \(\mathrm iA=G\)，
而且生成元定义域恰为 \(D(G)\)。
还须证明相同的零时导数确实决定整个群。
事实上，对任意具有同一生成元 \(G\) 的两个强连续酉群 \(U,V\)，
固定 \(t\in\mathbb R\)、\(x\in D(G)\)，令
\[
 F(s)=U(t-s)V(s)x\qquad(s\in\mathbb R).
\]
由定义域不变性，\(V(s)x\in D(G)\)。为核对对移动向量求导的步骤，写出
\[
 \begin{split}
 \frac{F(s+h)-F(s)}{h}
 &=U(t-s-h)\frac{V(s+h)x-V(s)x}{h}\\
 &\quad+\frac{U(t-s-h)-U(t-s)}{h}V(s)x.
 \end{split}
\]
第一项趋于 \(U(t-s)GV(s)x\)：
先用酉性控制差商与其极限的差，再对固定极限向量使用强连续性。
第二项的输入 \(V(s)x\) 已固定且属于 \(D(G)\)，
轨道导数公式给出其极限为 \(-U(t-s)GV(s)x\)。
因此 \(F'(s)=0\)。对任意固定 \(y\in H\)，
标量函数 \(\langle F(s),y\rangle\) 导数为零，故 \(F\) 是常值。
比较 \(s=0\) 与 \(s=t\) 得 \(U(t)x=V(t)x\)。
由于 \(D(G)\) 稠密，且两个算子有界，这个等式延伸到所有 \(x\in H\)。
应用于已构造的 \(V\)，就得到 \(U(t)=e^{\mathrm itA}\)。

若另一个自伴算子 \(B\) 也满足 \(U(t)=e^{\mathrm itB}\)，
\cref{b1:prop:29-spectral-unitary-group}的定义域结论给出
\(D(B)=D(G)\)、\(\mathrm iBx=Gx\)。故 \(B=A\)，包括定义域和作用都相同。
自伴算子到酉群的反方向已经由该命题完整证明。
\end{proof}

这个对应把方程
\[
 u'(t)=\mathrm iAu(t),\qquad u(0)=x\in D(A)
\]
的演化写成 \(u(t)=U(t)x\)，其范数保持为 \(\|x\|\)。
对任意 \(x\in H\)，连续轨道仍有定义；它具有强导数当且仅当 \(x\in D(A)\)。
由于 \(U(t)D(A)=D(A)\)，域外的初值在任何固定时间都不会进入该域。
若把指数约定改为 \(e^{-\mathrm itA}\)，相应的差商生成元便改为
\(-\mathrm iA\)；符号与定义域都属于对应关系的一部分。
